\documentclass[10pt, a4paper, twocolumn]{article}
\usepackage[utf8]{inputenc}
\usepackage[T1]{fontenc}
\usepackage{lmodern}
\usepackage[spanish]{babel}
\usepackage[margin=1.5cm]{geometry}
\usepackage{amsmath, amssymb}
\usepackage{graphicx}
\usepackage{booktabs}
\usepackage{caption}
\usepackage{tabularx}
\usepackage{microtype} % Soluciona la mayoría de los underfull/overfull de justificación
\usepackage{titlesec} % Configura los títulos para evitar overfulls y rupturas raras

% Evitar overfull/underfull en los títulos de sección en modo twocolumn
\titleformat{\section}{\large\bfseries\raggedright}{\thesection.}{0.5em}{}

\renewcommand{\abstractname}{Resumen}
\renewcommand{\figurename}{Fig.}
\renewcommand{\tablename}{Tabla}

\title{\Large \textbf{Plantilla de Curso con Organización Estilo IEEE - Hito 1\\ [0.2cm] TÍTULO DEL PROYECTO / HITO 1}}
\author{
    \textbf{Autores:} [Nombres de los estudiantes] \\
    \textbf{Carrera:} [Nombre de la carrera] \\
    \textbf{Asignatura:} IMT-221 — Circuitos Electrónicos III \\
    \textbf{Grupo:} [Número/Nombre del Grupo] \quad \textbf{Fecha:} [Fecha de entrega]
}
\date{}

\begin{document}
\maketitle

\begin{abstract}
[Un solo párrafo. Resuma el problema de ingeniería (protección ADC y supresión inductiva), el método (modelado, simulación, hardware), la evidencia principal validada (ej. error RMSE, voltaje de clamp medido) y la conclusión final. Manténgalo autocontenido y por debajo de 250 palabras.]
\end{abstract}

\textbf{\textit{Palabras clave—} Diodo Zener, Flyback, Protección Analógica, RMSE, LTspice, Python.}

\section{Introducción y Problema de Ingeniería}
Incluya: qué componente se está protegiendo (ADC), por qué existen sobretensiones o transitorios inductivos destructivos, el objetivo general del Hito 1 y los límites o rangos de voltaje relevantes. No coloque resultados finales aquí.

\section{Requisitos de Diseño y Supuestos}
Defina la topología empleada (Mecatrónica, Biomédica, contingencias de 3.3V), rango de entrada, y componentes.

\begin{table}[h]
\centering
\caption{Requisitos y Parámetros del Circuito}
\begin{tabularx}{\linewidth}{X r >{\raggedright\arraybackslash}X}
\toprule
\textbf{Parámetro} & \textbf{Condición} & \textbf{Justificación} \\
\midrule
$V_{in}$ & [-12, +24]\,V & Guía Hito 1 \\
$V_{ADC}$ perm. & [0, +5.1]\,V & Req. seguridad \\
$V_{ZT}$ Zener & 5.1\,V & Datasheet \\
$r_z$ (dinámica) & $15\,\Omega$ & Modelo PWL \\
$R_s$ (serie) & [ ]\,$\Omega$ & Cálculo + BOM \\
$R_{shunt}$ & $0.1\,\Omega$ & Anexo Hito 1 \\
\bottomrule
\end{tabularx}
\end{table}

\section{Modelo Analítico y Cálculos}
Muestre la deducción paso a paso. Plantee el modelo PWL del Zener:
\begin{equation}
    V_Z = V_{Z0} + r_z I_Z
\end{equation}
\begin{equation}
    V_{Z0} = V_{ZT} - I_{ZT} r_z
\end{equation}
Incluya el cálculo de los límites de $R_s$, indicando el peor caso térmico ($V_{in,max}=24\text{ V}$, $I_L=0$) y justificando la potencia disipada $P_Z$. Si existe discrepancia entre la guía, la BOM y su cálculo analítico, documéntela aquí y explique su decisión de ingeniería.

\section{Modelo Computacional en Python}
Describa el script utilizado para evaluar la respuesta teórica. Incluya las ecuaciones programadas, el barrido (\textit{sweep}), y cómo evaluó las regiones de operación del diodo.

\textit{[Insertar Fig. 1: Transferencia analítica $V_{out}$ vs $V_{in}$]}

\textit{[Insertar Fig. 2: Potencia del Zener $P_Z$ vs $V_{in}$]}

Explique en el texto qué demuestra cada figura generada por Python.

\section{Simulación en LTspice}
Documente el esquemático simulado. Indique los modelos de diodos utilizados, configuración de la simulación (\texttt{.dc}, \texttt{.tran}) y el impacto del diodo Flyback en la simulación de la carga inductiva.

\textit{[Insertar Fig. 3: Captura de resultados LTspice]}

Explique las diferencias encontradas respecto al modelo analítico simplificado.

\section{Implementación Física y Procedimiento Experimental}
Detalle el montaje en el laboratorio, los canales del osciloscopio, referencia de tierra común, y orientaciones físicas de los diodos (Zener, limitadores, y el diodo Flyback en paralelo a la carga).

\textit{[Insertar Fig. 4: Foto o esquemático del banco de pruebas]}

\section{Resultados}
Presente los datos extraídos objetivamente.

\begin{table}[h]
\centering
\caption{Comparación de Umbrales Críticos}
\begin{tabularx}{\linewidth}{X r r r c}
\toprule
\textbf{Magnitud} & \textbf{Teoría} & \textbf{LTspice} & \textbf{Físico} & \textbf{Ud.} \\
\midrule
Clamp sup. & [ ] & [ ] & [ ] & V \\
Clamp inf. & [ ] & [ ] & [ ] & V \\
\bottomrule
\end{tabularx}
\end{table}

\section{Comparación Cuantitativa y Validación}
Valide que las señales comparadas correspondan al mismo nodo, unidades y referencia. Reporte el error calculando el Root Mean Square Error (RMSE) tras una alineación temporal justificable:
\begin{equation}
    \text{RMSE} = \sqrt{\frac{1}{N}\sum_{k=1}^{N}(y_k-\hat{y}_k)^2}
\end{equation}

\textit{[Insertar Fig. 5: Superposición gráfica de CSV: Teoría vs LTspice vs Osciloscopio]}

Declare cualquier técnica de interpolación empleada.

\section{Discusión}
Interprete las discrepancias observadas en la Fig. 5. Analice el efecto de la resistencia dinámica $r_z$, tolerancias de los componentes, caídas directas no ideales, instrumentación, y limitaciones del modelo PWL de Python frente al hardware real.

\section{Conclusiones}
Responda concretamente si se logró la protección del ADC, basándose en la evidencia. Mencione qué limitación persiste en el diseño actual y qué aprendizaje analítico o práctico se lleva hacia el Hito 2 (sensado diferencial BJT).

\section*{Referencias}
[1] Guía del Hito 1 y Anexos, IMT-221, UCB Tarija. \\
\relax [2] Datasheet BZX55C5V1, [Fabricante]. \\
\relax [3] Jaeger, Blalock \& Blalock, \textit{Microelectronic Circuit Design}, 6th ed.

\section*{Apéndice A - Auditoría de Uso de IA}
(Requisito obligatorio de la asignatura). Describa:
1) Herramientas utilizadas.
2) Prompt o interacción relevante para el desarrollo.
3) Identificación de al menos un error, simplificación o "alucinación" del modelo detectada por el equipo, y cómo fue corregida.

\end{document}
